\section{The Holonomy Interface}
\label{p13:sec:holonomy}

\noindent
The mechanism of \xref{p13:sec:mechanism} was given in the vocabulary of prediction and precision. This section gives the same mechanism a second reading, in the geometric vocabulary the series has developed elsewhere\,---\,the vocabulary of holonomy, of the phase a system accumulates in traversing a closed path. The second reading is not a rival to the first but its complement: where the predictive account describes what the receiving mind computes, the geometric account describes the \emph{shape} of the process by which a promise, held across time, either sustains the trust placed in it or drains it away. The two readings will be seen to converge, and their convergence prepares two later moments of the paper: the quantum-dynamical structure of relational trust (\xref{p13:sec:relational}), for the geometric phase is, in its origin, a concept of quantum dynamics; and the inversion of value by which the non-binding vow proves the purer instrument (\xref{p13:sec:praxis}), for that inversion turns on exactly the distinction this section now draws.

A discipline must be stated at the outset, as it was for the mechanism. The geometric phase is employed here as a \emph{structural analogy}, developed to the point of usefulness but not to full formal closure; the formal programme that would furnish the phase with a rigorous carrier, a definite fibre and connection, is carried by another paper of the series, to which the present section forward-refers rather than reproducing it. What is claimed here is that the structure of a promise sustained across time is the structure of a system transported around a closed loop, and that the distinction between a promise that generates trust and one that drains it is the distinction between a loop that accumulates a positive holonomy and one that accumulates none or a negative one. The claim is structural; its full formalisation lies elsewhere.

\subsection{Trust as a quantity transported around a loop}
\label{p13:sub:holo-loop}

Consider a promise not at the instant of its making but across the time of its keeping. A relation in which a vow has been made returns, again and again, to situations that resemble the situation of the vow: occasions on which the commitment is tested, recalled, re-enacted, or merely lived through\,---\,the daily, monthly, yearly recurrences in which what was promised is, or is not, borne out. The relation thus traces, across its time, a closed loop in a space of situations: it departs from the moment of the vow, passes through the varied circumstances of a shared life, and returns to situations that are, in the relevant respect, the same situation again. The question the geometric reading asks is what the relation \emph{brings back} when it returns. Does it return to the situation of the vow unchanged, having merely gone around and come back? Or does it return having accumulated something\,---\,a deepening, a surplus, a phase that was not there at the start?

This is the question of holonomy. A system transported around a closed loop in a space with the right structure does not, in general, return to its starting state; it returns rotated, shifted, carrying a phase that records the path it took\,---\,the geometric phase, the anholonomy, of which the Berry phase and the Pancharatnam phase are the canonical instances~\parencite{berry1984,pancharatnam1956}. The phase is not a function of where the system ended\,---\,it ended where it began\,---\,but of the loop it traversed to get back: it is the trace, in the returned state, of the journey. And the sign and magnitude of that phase is exactly the distinction the series has used to tell a good cycle from a bad. A loop that returns its system unchanged has zero holonomy: it has gone around and come back to the same place with nothing to show for the going, the mere repetition the series has called the bad infinity. A loop that returns its system enriched, carrying a positive phase, has positive holonomy: it has returned to its root and yet come back as more, the spiral rather than the circle. And a loop that returns its system depleted, having leaked phase rather than gained it, has negative holonomy: the cycle of net extraction, the relation that winds down. $\holo$ of the loop, the holonomy of the relation's traversal, is the geometric measure of whether the promise, lived across time, deepens, repeats, or drains.

\subsection{The non-coercive promise as positive, adiabatic holonomy}
\label{p13:sub:holo-adiabatic}

The non-coercive promise, kept, is the loop of positive holonomy, and it is so in the specific manner the geometric phase requires: \emph{adiabatically}. The geometric phase is the phase accumulated by a system carried around its loop slowly enough that it is not driven by the transport but follows it, exchanging no energy with its surroundings, its evolution governed from within rather than pumped from without. This adiabatic condition is the geometric image of exactly what \xref{p13:sec:freedom} called the inferential cause and \xref{p13:sec:mechanism} called the self-committing and wagered channels. The promise kept without a sword is kept adiabatically: its loop is closed not by an external compulsion driving the parties around it but by the parties' own commitment, sustained from within, returning upon itself by the completeness of what each freely brings to it. Nothing external pumps the cycle; the parties are not driven around the loop by a threatened consequence; the closure is endogenous, and the phase it accumulates is, for that reason, a true geometric phase\,---\,a deepening generated by the loop's own internal structure, the surplus the series has identified with the good cycle. The non-coercive promise is the adiabatic loop of positive holonomy: trust sustained and deepened by the relation's own internal completeness, requiring no external pump to hold it closed.

This gives the geometric content of the claim of \xref{p13:sec:freedom}, that the promise without force is the proper case and not the defective one. In the geometric vocabulary, the want of an external pump is not a deficiency of the loop but the condition of its phase being geometric at all. A phase that is pumped from outside is not a geometric phase; it is a dynamical phase, driven by the energy fed in, and it is not the trace of the loop's own structure but of the driving. The adiabatic loop\,---\,the one no external force drives\,---\,is the only one whose accumulated phase is purely the loop's own; and so the promise that relies on no sword is the only one whose deepening is genuinely its own, generated from within rather than supplied from without. The geometric reading thus recovers, in its own terms, the result that the absence of force is the condition of the promise's deepening being the promise's own.

\subsection{Debt and the coerced contract as the zero-holonomy line}
\label{p13:sub:holo-debt}

The contrast that throws this into relief is the coerced contract, and behind it the structure of debt, which the geometric reading exposes as something quite different from a loop. A debt is not a cycle that returns upon itself and deepens; it is a line, drawn from a present that is mortgaged to a future that must discharge it. The debtor does not traverse a loop that brings back a surplus; he is held, by the threat of the sanction, on a track that runs from the incurring of the obligation to its forced discharge, and the track is held straight\,---\,prevented from closing into any loop of its own\,---\,by the external force that drives it. This is the geometric image of the coerced contract and of debt as its limit: a path held open and directed by an external pump, accumulating no geometric phase because it is not permitted to close upon itself, its motion not the endogenous traversal of a loop but the driven progress along a line.

Where such a path is forced into the appearance of a cycle\,---\,where a coerced relation is made to repeat\,---\,its holonomy is zero or negative, never the positive phase of the adiabatic loop. Zero, in the case of mere forced repetition: the relation goes around and comes back to the same place, the obligation discharged and re-incurred, with nothing brought back, because the closure was imposed from outside and the loop accumulated no phase of its own\,---\,the bad infinity of a cycle held shut by force. Negative, in the case of net extraction: where the external pump drives the loop so as to draw phase out of it, the relation winds down, depleted turn by turn, the structure the series has identified with the parasitic and the exploitative. The coerced relation cannot accumulate positive holonomy, because positive holonomy is the signature of an adiabatic, endogenously closed loop, and the coerced relation is precisely the one that is not adiabatic: it is pumped from without, and a pumped loop's phase is not its own. The sword, on this reading, does not merely fail to deepen the bond; it forecloses the deepening, by substituting a driven line or a forced repetition for the free, adiabatic loop that alone could deepen.

\subsection{The spectral gap and the \zh{甘え} problem}
\label{p13:sub:holo-amai}

There is a further structure the geometric reading illuminates, which the series left open in its treatment of \zh{甘え}\,---\,of the sweet, regressive presuming-upon-another's-indulgence whose ambivalence that discussion could mark but not resolve. The unresolved question there was where, in the space of a relation's possible motions, the line falls between a presuming-upon-another that the relation can absorb and metabolise, returning it as deepened intimacy, and a presuming-upon-another that the relation cannot absorb, which accumulates as deviation and drives the cycle toward divergence. The geometric reading offers the missing structure, in the form of a spectral gap. A loop's capacity to absorb a perturbation and return adiabatically to its closed traversal\,---\,rather than being driven by the perturbation onto a divergent path\,---\,depends on a gap in the spectrum of the dynamics governing the loop: a perturbation slower than the gap is followed adiabatically and metabolised into the loop's own phase, while a perturbation faster than the gap drives the system across the gap, out of the adiabatic regime, onto a trajectory that no longer closes. The benign \zh{甘え} is the perturbation the loop can follow adiabatically, absorbed into a deepened holonomy; the corrosive \zh{甘え} is the perturbation that exceeds the gap, driving the relation non-adiabatically toward the accumulation of deviation and divergence that \xref{p13:sub:mech-forged} described in the predictive register as the collapse of a forged precision. The spectral gap is the structure that was missing; it is offered here as the geometric correlate of the relation's capacity to bear what is presumed upon it, and the point at which the present paper repays, in part, a debt the earlier discussion left standing.

\subsection{The convergence of the two readings}
\label{p13:sub:holo-convergence}

The geometric reading and the predictive reading of \xref{p13:sec:mechanism} are, in the end, one account seen twice, and their convergence is worth stating plainly, for it is what licenses the use of both. The adiabatic, endogenously closed loop of positive holonomy is the geometric portrait of the relation whose trust is sustained by the three coupled channels from within, requiring no external pump\,---\,the genuine case. The driven line, and the forced loop of zero or negative holonomy, is the geometric portrait of the relation whose precision is injected from without rather than earned by a real likelihood\,---\,the coerced case and the forged case alike. What the predictive reading describes as a precision earned by a real self-committing likelihood versus a precision injected by external pressure, the geometric reading describes as a phase accumulated by an adiabatic loop versus a phase pumped into a driven path; and these are the same distinction, the distinction between a process governed from within and one driven from without, rendered in two vocabularies. That the two independent readings converge on the same line\,---\,between the endogenous and the pumped, the adiabatic and the driven, the genuine and the forged\,---\,is the strongest internal evidence the paper can offer that the line is real and not an artefact of either vocabulary.

This convergence carries forward in two directions the paper will take up. It points toward \xref{p13:sec:relational}, where the recognition that the geometric phase is, in its origin, a concept of quantum dynamics\,---\,the Berry phase being the phase of an adiabatically transported quantum state\,---\,allows the geometric language used here to be seen for what it has been all along, a borrowing from the dynamics of which relational trust will there be shown to share the structure. And it points toward \xref{p13:sec:praxis}, where the distinction between the adiabatic loop, whose phase is its own, and the pumped path, whose phase is borrowed, becomes the geometric ground of the inversion by which the non-binding vow\,---\,the loop that relies on no external pump\,---\,proves the purer realisation of what a promise is, and the enforceable instrument, with its admixture of the external pump, the diluted one.
